% ============================================================
% WearLens report -- Sections 2-5 (Overview, Data, Users & Challenges, Tasks)
% Paste into the course LaTeX template. Requires \usepackage{booktabs} (and {tabularx} for Table 3).
% All numbers are traceable to the repository: data/eda_tool_life.csv, data/eda_spearman.csv,
% data/features_enriched.csv, docs/07_verification.md.
% Lines marked TODO need a team decision or a citation you have checked yourselves.
% ============================================================

\section{Overview}
End mills wear gradually, and a worn tool degrades surface quality and process stability.
Shops therefore usually replace tools after a fixed period, which wastes much of the remaining life
of most tools~\cite{denkena2023}. Measuring wear directly requires stopping the machine and inspecting
the tool under a microscope, so engineers would like to monitor wear through signals the machine already
records. This is difficult on a \emph{fleet} of machines: the same tool model lasts noticeably longer on one
machine than on another, the same signal behaves differently across machines, and many signal changes are
not caused by wear at all.
\textbf{WearLens} is a visual analytics system for process engineers who introduce tool-condition
monitoring on several machine tools. It helps them
(i) choose replacement times that do not waste tool life,
(ii) find a machine-control signal that tracks wear consistently across machines, and
(iii) tell real wear from signal changes with other causes.

\section{Data and Data Processing}
\paragraph{Source.}
We use the open dataset of Denkena \etal~\cite{denkena2023}
(Mendeley Data, \url{https://doi.org/10.17632/zpxs87bjt8}, CC~BY~4.0).
Nine solid-carbide end mills of the same model (D\,8\,mm, 4 edges) were each worn from new to
the end of life ($\mathrm{VB}\approx150\,\mu$m) by shoulder milling cast iron on three 5-axis machines
(M1: T1--T3, M2: T4--T6, M3: T7--T9), using identical NC code and process parameters.
Each \emph{run} is one straight milling pass of about 4\,s cutting, stored as one HDF5 file;
44 runs form one 2\,mm layer. The dataset holds \textbf{6,418 runs} (10.6\,GB).

\paragraph{Attributes.}
Table~\ref{tab:data} lists the attributes we use. Each run contains multivariate time series from two
independent measurement chains: the machine control (500\,Hz) and a Kistler dynamometer below the
workpiece (25\,kHz). In our scenario, only the machine-control signals are available in production;
the dynamometer serves as a reference from the trial (Section~\ref{sec:users}).

\begin{table}[t]
\centering\small
\caption{Attributes used by WearLens (per run unless noted).}
\label{tab:data}
\begin{tabular}{@{}lll@{}}
\toprule
Attribute & Type & Description \\
\midrule
Machine M, tool T & categorical & 3 machines, 9 tools (3 per machine) \\
Run R & ordinal & pass counter per tool (removed runs not counted) \\
Contact time C & quantitative [s] & cumulated cutting time before the run \\
Flank wear VB & quantitative [$\mu$m] & mean of 4 edges, microscope \\
Layer, pass position $x$ & ordinal / quant.\ [mm] & derived from \texttt{tool\_position\_x} \\
Machine signals & time series, 500\,Hz & spindle torque, feed-axis torque (M1) or force (M2/M3), \\
 & & position control deviation, tool position \\
Dynamometer forces & time series, 25\,kHz & $F_x, F_y, F_z$ (reference only) \\
\bottomrule
\end{tabular}
\end{table}

\paragraph{Preprocessing (Python, offline).}
We verified every file against its SHA-256 checksum and checked that the labels inside each file
match \texttt{filelist.csv} (6,418/6,418).
For each run we
(1) detect the cutting window from the dynamometer's $xy$ resultant force
(50\,ms moving RMS above 30\,\% of its 95th percentile; detected length 3.91--4.05\,s, consistent with
$\approx$4\,s of contact time per run),
(2) trim 0.15\,s at both ends to exclude tool entry/exit and the up-to-100\,ms misalignment between the two
measurement chains, and
(3) compute 21 features in this window: 8 from the machine control
(e.g., mean $|$spindle torque$|$, mean $|$feed-axis effort$|$, standard deviation of the position control
deviation) and 13 from the dynamometer (e.g., mean resultant force, amplitude at the tooth-passing
frequency of 530\,Hz).
We take absolute values of the spindle torque because its sign convention differs between machines.
We also derive the layer index and the pass position within a layer, and store a 200-bin envelope of three
signals per run for the drill-down view.
The web application loads only compact files (a 1.2\,MB CSV with one row per run, 7.7\,MB of envelopes, and
0.7\,MB of raw 20\,ms force snippets). Normalization, correlation, tool-life estimation and layer
diagnostics are computed in the browser so that users can change their parameters
(Section~7). % TODO: adjust the section number to your template
To confirm the numbers, we recomputed the features independently for 40 random runs (maximum relative error
$4\times10^{-4}$) and checked the values shown in the interface against the Python results.

\paragraph{Known data limitations.}
We show the following limitations in the interface and take them into account in our conclusions:
VB was measured only every one or two layers ($\approx$40 runs) and is linearly interpolated in between;
the first runs of T8 are missing (data start at $\mathrm{VB}=34\,\mu$m);
spindle torque and feed-axis force on M2 may be aliased;
the coordinate system of T7 and T8 was rotated by about 1--2$^\circ$;
and all M2 tools come from a different tool batch than the others, so the effect of the machine cannot be
separated from that of the batch.

\section{Target Users and Analysis Challenges}
\label{sec:users}
\paragraph{Target users.}
Our users are \textbf{process or maintenance engineers who introduce tool-condition monitoring on a fleet
of machine tools}. Before deployment, such engineers run trials like the one in this dataset: tools are
worn on several machines while machine signals are recorded, and wear is checked with a microscope;
a dynamometer may be added in the trial. In production, however, dynamometers are rarely available
because they are expensive, reduce the stiffness of the setup, and constrain the workpiece.
% TODO: support this sentence with a tool-condition-monitoring review you have checked yourselves.
Only the signals of the machine control remain. The engineers know machining well but have no time to
build and validate custom analysis scripts for every machine.

\paragraph{CH1: When should a tool be replaced?}
Tool life varies across machines even with identical tools and parameters.
To reach $\mathrm{VB}=140\,\mu$m, the M2 tools need 81.5--84.0\,min of contact time, whereas the M1 and M3
tools need only 53.5--57.9\,min, and almost all of this difference arises in the steady middle phase of
wear. A single fixed replacement time that is safe for every tool must follow the shortest-lived tool;
this leaves 34--36\,\% of the life of each M2 tool unused. The threshold itself is a trade-off between
quality risk and tool cost, which only the engineer can decide.

\paragraph{CH2: Which production signal can serve as a wear indicator across machines?}
There are many candidate features, and each can be normalized in several ways. Raw values cannot be
compared across machines: pooled over all tools, the Spearman correlation between mean spindle torque and
VB is only 0.57, but 0.97 after normalizing within each machine.
Some features even change in opposite directions on different machines: the 530\,Hz tooth-passing
amplitude rises with wear on M1 and M2 but falls on M3.
Correlations within a single tool are also optimistic, because VB is mostly interpolated from contact time,
so any monotonic drift correlates with it.
Judging which indicator is trustworthy therefore requires comparing curves across tools and machines,
not reading off a single coefficient.

\paragraph{CH3: Is a signal change really wear?}
Signals also change for reasons other than wear. In tool T1, every signal in layer 15 jumps by about
20\,\% and drops again in the next layer. In tool T9, the force rises by 26--52\,\% within each of layers
10--14 and resets at every new layer. Both effects appear in the machine signals and in the dynamometer,
two independent measurement chains, so they are unlikely to be sensor faults. Wear, however, cannot
decrease. An alarm based on the indicator alone would therefore suggest that wear suddenly accelerated
and trigger an early replacement. Finding such cases requires looking at the raw signals and relating them
to the layer structure, and their causes (for example, the machine stopping between layers for
measurement) are not recorded in the data.

\paragraph{Why visual analytics?}
None of these challenges can be answered by a simple query or a single static chart.
The answers depend on combinations of machine, tool, feature, normalization and threshold, and the
interesting cases (such as T1 and T9) were not known before we explored the data.
A fully automatic model is also not appropriate. With only nine tools (three per machine), it cannot be
validated for new machines; most VB labels are interpolated; and unusual segments need a domain
explanation. Most importantly, engineers must \emph{see the evidence} before they trust an indicator on the
shop floor. WearLens therefore keeps the engineer in the loop: it provides the derived measures and lets
the engineer judge them.

\section{Low-Level Tasks}
We derived seven tasks from the challenges (Table~\ref{tab:tasks}).
T1--T2 address replacement timing (CH1),
T3--T4 the choice of an indicator (CH2),
T5--T6 the separation of wear from other effects (CH3), and
T7 connects CH2 back to CH1 by translating an indicator level into a replacement rule
that needs no microscope.

\begin{table}[t]
\centering\small
\caption{Low-level tasks and the challenges they come from.}
\label{tab:tasks}
\begin{tabularx}{\linewidth}{@{}lXll@{}}
\toprule
ID & Task & Type & From \\
\midrule
T1 & Overview the wear curves of all tools and compare in which wear phase the machines differ & overview, compare & CH1 \\
T2 & For a chosen VB threshold, identify when each tool reaches it and derive how much life a fixed replacement time leaves unused & identify, derive & CH1 \\
T3 & Rank features by their correlation with wear for each tool and for all tools, under a chosen normalization and time range & rank, correlate & CH2 \\
T4 & Compare the normalized trend of one feature across the three machines to judge whether it transfers & compare & CH2 \\
T5 & Detect runs or layers that deviate from the wear trend and locate them in time and in the layer structure & detect outliers, locate & CH3 \\
T6 & Inspect the signals of selected runs and check whether machine signals and the dynamometer reference agree & details on demand, verify & CH3 \\
T7 & Read off the indicator level at which tools reach the VB threshold & correlate, derive & CH1+CH2 \\
\bottomrule
\end{tabularx}
\end{table}

% ---------------- bibliography entry (add to your .bib) ----------------
% @article{denkena2023,
%   author  = {Denkena, Berend and Klemme, Heinrich and Stiehl, Tobias H.},
%   title   = {Multivariate time series data of milling processes with varying tool wear and machine tools},
%   journal = {Data in Brief}, volume = {50}, pages = {109574}, year = {2023},
%   doi     = {10.1016/j.dib.2023.109574} }
% If your template has no \etal macro, replace "Denkena \etal" with "Denkena et al.".
